\chapter{Constrained Optimization: KKT Conditions and Second-Order Tests}

\graphicspath{{chapters/07_constrained_optimization/figures}}

\begin{abstract}
This chapter develops local optimality conditions for constrained problems:
feasible directions, tangent cones, constraint qualifications, first-order KKT
conditions, second-order tests, and representative examples.
\end{abstract}

\tikzset{
  ch7-axis/.style={-latex, line width=0.55pt, draw=slate!90},
  ch7-boundary/.style={draw=deepblue, line width=0.9pt},
  ch7-feasible/.style={draw=teal!80!black, fill=teal!8, line width=0.7pt},
  ch7-arrow/.style={-latex, line width=0.9pt},
  ch7-normal/.style={ch7-arrow, draw=deepblue},
  ch7-grad/.style={ch7-arrow, draw=crimson!85!black},
  ch7-tangent/.style={ch7-arrow, draw=teal!80!black},
  ch7-label/.style={font=\footnotesize, text=slate!95},
  ch7-small/.style={font=\scriptsize, text=slate!92},
  ch7-point/.style={circle, inner sep=1.3pt, fill=black, draw=black},
}

\section{Introduction}
Chapter~1 explained first- and second-order optimality conditions for
unconstrained problems, and Chapter~4 introduced the convex sets and convex
constraint classes that appear in constrained models. The new issue here is
that, near a feasible point, not every direction is admissible. Constrained
optimization therefore starts with a local question: how can one recognize or
certify a minimizer when the feasible set blocks most directions?



\section{Problem statement}
\label{sec:ch7-problem}
We begin with the standard constrained problem
\begin{equation}
  \label{eq:ch7-problem}
  \begin{aligned}
    \min_{\mathbf{x}} \quad & f(\mathbf{x}) \\
    \text{subject to}\quad & g_i(\mathbf{x}) \le 0, \qquad i=1,\dots,m, \\
    & h_j(\mathbf{x}) = 0, \qquad j=1,\dots,p,
  \end{aligned}
\end{equation}
where
\[
  f, g_i, h_j \in C^1.
\]

\subsection{Basic definitions}
The basic objects attached to \eqref{eq:ch7-problem} are the feasible set, the
active constraints, and the Lagrangian.

\begin{definition}[Feasible set]
We define the common domain by
\[
  D
  := \operatorname{dom} f
  \cap \bigcap_{i=1}^m \operatorname{dom} g_i
  \cap \bigcap_{j=1}^p \operatorname{dom} h_j
\]
and the feasible set by
\[
  F
  := \left\{\mathbf{x}\in D:
  g_i(\mathbf{x})\le 0 \ \forall i,\quad
  h_j(\mathbf{x})=0 \ \forall j\right\}.
\]
\end{definition}

\begin{definition}[Active constraints]
Following the notation of this chapter, an inequality constraint is \emph{active} (or
\emph{binding}) at $\mathbf{x}$ if it holds with equality. We denote the set of
active constraints at $\mathbf{x}$ by
\[
  \operatorname{Active}(\mathbf{x})
  :=
  \left\{g_i : g_i(\mathbf{x})=0,\ i=1,\dots,m\right\}
  \cup
  \left\{h_j : j=1,\dots,p\right\}.
\]
Thus the active constraints are precisely the active inequalities together with
all equality constraints. When active indices are meant, we use the same
notation $\operatorname{Active}(\mathbf{x})$.
\end{definition}

\begin{definition}[Lagrangian]
The associated Lagrangian is
\begin{equation}
  \label{eq:ch7-lagrangian}
  L(\mathbf{x},\boldsymbol{\lambda},\boldsymbol{\mu})
  := f(\mathbf{x})
  + \sum_{i=1}^m \lambda_i g_i(\mathbf{x})
  + \sum_{j=1}^p \mu_j h_j(\mathbf{x}).
\end{equation}
\end{definition}

\subsection{Equivalent descriptions of the same feasible set}
Before stating KKT, one modeling issue matters: the feasible set alone does not
determine the form of the optimality conditions. What also matters is how that
set is represented through equalities and inequalities. The same feasible set
may admit different descriptions. For example,
\[
  h(\mathbf{x}) = 0
  \iff \widehat h(\mathbf{x}) := h(\mathbf{x})^2 = 0
\]
and also
\[
  h(\mathbf{x}) = 0
  \iff
  \begin{cases}
    g_1(\mathbf{x}) := h(\mathbf{x}) \le 0, \\
    g_2(\mathbf{x}) := -h(\mathbf{x}) \le 0.
  \end{cases}
\]
This freedom is useful, but it is also dangerous: some reformulations destroy
the constraint qualifications needed for KKT. We will see an explicit example of
this failure later in the chapter.

So, KKT is attached to a
\emph{constraint representation}, not only to the feasible set itself. Two
formulations may define the same set $F$ and still behave differently from the
viewpoint of regularity, because they produce different gradients for the
active constraints. When we apply KKT, we must therefore work with a
representation whose constraints satisfy the required qualification. Some
constrained problems can also be reduced to unconstrained ones; in this
chapter, however, we keep the constrained formulation and derive the optimality
conditions directly from that representation.

\section{First-order KKT conditions}
\label{sec:ch7-kkt-first-order}
\begin{theorem}[First-order KKT conditions]
\label{thm:ch7-kkt-first}
Assume that $f,g_i,h_j\in C^1$, let $\mathbf{x}^\star$ be a local minimizer of
\eqref{eq:ch7-problem}, and assume that a constraint qualification holds at
$\mathbf{x}^\star$ (the relevant constraint qualifications are listed in
\Cref{sec:ch7-cq}). Then \(\exists\, \boldsymbol{\lambda}^\star \in \mathbb{R}^m\)
and \(\exists\, \boldsymbol{\mu}^\star \in \mathbb{R}^p\) such that
\begin{enumerate}
  \item
  \begin{equation}
    \label{eq:ch7-kkt-stationarity}
    \nabla_{\mathbf{x}} L(\mathbf{x}^\star,\boldsymbol{\lambda}^\star,\boldsymbol{\mu}^\star)
    =
    \nabla f(\mathbf{x}^\star)
    + \sum_{i=1}^m \lambda_i^\star \nabla g_i(\mathbf{x}^\star)
    + \sum_{j=1}^p \mu_j^\star \nabla h_j(\mathbf{x}^\star)
    = \mathbf{0};
  \end{equation}
  \item $\mathbf{x}^\star \in F$;
  \item $\lambda_i^\star \ge 0 \ \forall i$;
  \item $\lambda_i^\star g_i(\mathbf{x}^\star)=0 \ \forall i$.
\end{enumerate}
Condition \eqref{eq:ch7-kkt-stationarity} is the stationarity of the
Lagrangian, the second condition is primal feasibility, the third is dual
feasibility, and the fourth condition is complementary slackness.
\end{theorem}

The theorem is the main practical message. The rest of this section derives it
from the geometry of feasible directions, the role of constraint qualifications,
and the Farkas alternative.

\subsection{Feasible curves and the tangent cone}
\label{sec:ch7-tangent}
Suppose $\mathbf{x}^\star\in F$. If $\mathbf{x}^\star$ is an interior point of $F$,
then locally the problem behaves like an unconstrained one, so we should again
expect $\nabla f(\mathbf{x}^\star)=\mathbf{0}$. The interesting case is
$\mathbf{x}^\star\in \partial F$, where only some directions remain feasible.

To make that precise, consider all smooth feasible curves issuing
from $\mathbf{x}^\star$:
\[
  S(\mathbf{x}^\star)
  := \left\{
    \mathbf{x}(\cdot)\in C^1:
    \mathbf{x}(0)=\mathbf{x}^\star,\ 
    \mathbf{x}(t)\in F \text{ for } t\in[0,t_{\max}]
  \right\}.
\]

\begin{definition}[Tangent cone]
\label{def:ch7-tangent-cone}
The tangent cone to $F$ at $\mathbf{x}^\star$ is
\[
  T_F(\mathbf{x}^\star)
  := \left\{
    \mathbf{d}\in\mathbb{R}^n:
    \exists\, \mathbf{x}(\cdot)\in S(\mathbf{x}^\star)
    \text{ with }
    \dot{\mathbf{x}}(0)=\mathbf{d}
  \right\}.
\]
\end{definition}

\begin{figure}[ht]
    \centering
    \incfig{tangent-cone}
    \caption{Smooth feasible curves $\mathbf{x}_1(t)$ and $\mathbf{x}_2(t)$
    issuing from $\mathbf{x}^\star\in F$. Their initial velocities
    $\mathbf{d}_1=\dot{\mathbf{x}}_1(0)$ and $\mathbf{d}_2=\dot{\mathbf{x}}_2(0)$
    are feasible first-order directions, hence they belong to the tangent cone
    $T_F(\mathbf{x}^\star)$.}
    \label{fig:ch7-tangent-cone}
\end{figure}
This is the collection of all limiting directions of feasible trajectories
passing through $\mathbf{x}^\star$; see \Cref{fig:ch7-tangent-cone}.

\begin{claim}[Tangent-cone necessary condition]
\label{clm:ch7-tangent-necessary}
If $\mathbf{x}^\star\in F$ is a local minimizer of \eqref{eq:ch7-problem}, then
\begin{equation}
  \label{eq:ch7-tangent-necessary}
  \nabla f(\mathbf{x}^\star)^\top \mathbf{d} \ge 0
  \qquad \forall \mathbf{d}\in T_F(\mathbf{x}^\star).
\end{equation}
\end{claim}
\begin{proof}
Take any feasible curve $\mathbf{x}(\cdot)\in S(\mathbf{x}^\star)$. Since
$\mathbf{x}^\star$ is a local minimizer, the one-variable function
$t\mapsto f(\mathbf{x}(t))$ has a local minimum at $t=0$, hence
\[
  \left.\frac{d}{dt} f(\mathbf{x}(t))\right|_{t=0}\ge 0.
\]
By the chain rule,
\[
  \left.\frac{d}{dt} f(\mathbf{x}(t))\right|_{t=0}
  = \nabla f(\mathbf{x}^\star)^\top \dot{\mathbf{x}}(0).
\]
Now set $\mathbf{d}=\dot{\mathbf{x}}(0)$. By definition,
$\mathbf{d}\in T_F(\mathbf{x}^\star)$, and \eqref{eq:ch7-tangent-necessary}
follows.
\end{proof}

The meaning of \eqref{eq:ch7-tangent-necessary} is straightforward: along
every feasible first-order direction, the objective must not decrease.

\begin{figure}[ht]
    \centering
    \incfig{kkt-condition}
    \caption[Geometry of equality and inequality KKT conditions]{Geometry of the
    first-order KKT conditions in the two basic cases. The gradients
    $\nabla h(\mathbf{x}^\star)$ and $\nabla g(\mathbf{x}^\star)$ are normal to
    their level sets because, along a level set, the constraint is constant and 
    therefore the directional derivative is
    zero. Left: for $h(\mathbf{x})=0$, feasible directions satisfy
    $\nabla h(\mathbf{x}^\star)^\top\mathbf{d}=0$, so
    $T_F(\mathbf{x}^\star)$ is the tangent space. Since both $\mathbf{d}$ and
    $-\mathbf{d}$ are feasible, first-order optimality forces
    $\nabla f(\mathbf{x}^\star)$ to have no tangent component:
    $\nabla f(\mathbf{x}^\star)=-\mu\nabla h(\mathbf{x}^\star)$, with
    $\mu\in\mathbb{R}$ free in sign. Right: for $g(\mathbf{x})\le0$,
    $\nabla g(\mathbf{x}^\star)$ points toward increasing $g$, hence outward
    from the feasible side $g<0$. Feasible directions satisfy
    $\nabla g(\mathbf{x}^\star)^\top\mathbf{d}\le0$, and the condition
    $\nabla f(\mathbf{x}^\star)^\top\mathbf{d}\ge0 \ \forall
    \mathbf{d}\in T_F(\mathbf{x}^\star)$ forces
    $\nabla f(\mathbf{x}^\star)=-\lambda\nabla g(\mathbf{x}^\star)$ with
    $\lambda\ge0$.}
    \label{fig:ch7-kkt-condition}
\end{figure}

If $\mathbf{x}^\star\in \operatorname{int} F$, then
$T_F(\mathbf{x}^\star)=\mathbb{R}^n$, so \eqref{eq:ch7-tangent-necessary} reduces
to the unconstrained condition
\[
  \nabla f(\mathbf{x}^\star)=\mathbf{0}.
\]
So far, however, \eqref{eq:ch7-tangent-necessary} is not very convenient to
check directly. The purpose of KKT is to rewrite it in terms of the gradients
of the constraint functions themselves.

\subsection{Case 1: one equality constraint}
\label{subsec:ch7-one-equality}
Consider first the problem
\[
  f(\mathbf{x}) \to \min_{\mathbf{x}},
  \qquad
  h(\mathbf{x})=0.
\]
At a feasible point $\mathbf{x}^\star$ with $\nabla h(\mathbf{x}^\star)\neq \mathbf{0}$,
the tangent cone is
\begin{equation}
  \label{eq:ch7-tangent-equality}
  T_F(\mathbf{x}^\star)
  = \left\{
    \mathbf{d}\in\mathbb{R}^n:
    \nabla h(\mathbf{x}^\star)^\top \mathbf{d}=0
  \right\}.
\end{equation}
This geometry is shown in the left panel of \Cref{fig:ch7-kkt-condition}.
Indeed, first-order feasibility means that along every feasible curve the
constraint stays constant, so its derivative at $t=0$ must vanish.

Now apply \eqref{eq:ch7-tangent-necessary}. Since
$\mathbf{d}\in T_F(\mathbf{x}^\star)$ implies $-\mathbf{d}\in T_F(\mathbf{x}^\star)$,
the inequality
\[
  \nabla f(\mathbf{x}^\star)^\top \mathbf{d}\ge 0
  \qquad \forall \mathbf{d}\in T_F(\mathbf{x}^\star)
\]
forces equality:
\[
  \nabla f(\mathbf{x}^\star)^\top \mathbf{d}=0
  \qquad \forall \mathbf{d}\in T_F(\mathbf{x}^\star).
\]
So $\nabla f(\mathbf{x}^\star)$ is orthogonal to the tangent space
\eqref{eq:ch7-tangent-equality}. The normal space is
spanned by $\nabla h(\mathbf{x}^\star)$, hence $\exists\, \mu\in\mathbb{R}$ such
that
\[
  \nabla f(\mathbf{x}^\star) + \mu \nabla h(\mathbf{x}^\star)=\mathbf{0}.
\]
This is exactly the equality-constrained form of KKT stationarity.

Two edge cases are worth keeping in mind.
\begin{itemize}
  \item If $\nabla f(\mathbf{x}^\star)=\mathbf{0}$, then the condition collapses to
  the unconstrained one from Chapter~1.
  \item If $\nabla h(\mathbf{x}^\star)=\mathbf{0}$, then the linearized
  constraint becomes trivial and imposes no restriction on the direction
  $\mathbf{d}$. As a consequence, the first-order approximation suggests that
  every direction is feasible, whereas the actual feasible set may admit only a
  much smaller set of directions. This is the first sign that a constraint
  qualification is needed.
\end{itemize}

\subsection{Constraint qualifications}
\label{sec:ch7-cq}
Constraint qualifications are introduced precisely to rule out such
degenerate situations. Informally, a constraint qualification says that near
the current point, the local feasible geometry is captured correctly by the
linearized active constraints, namely the constraints in
$\operatorname{Active}(\mathbf{x})$.

\begin{definition}[Common constraint qualifications]
\label{def:ch7-cq}
At a feasible point $\mathbf{x}^\star\in F$, each of the following is a standard
constraint qualification.
\begin{enumerate}
  \item \textbf{LCQ (linear constraint qualification).}
  All constraint functions are affine with nonzero linear part.

  \item \textbf{LICQ (linear independence constraint qualification).}
  The family
  \[
    \left\{\nabla g_i(\mathbf{x}^\star): i\in\operatorname{Active}(\mathbf{x}^\star)\right\}
    \cup
    \left\{\nabla h_j(\mathbf{x}^\star): j\in\operatorname{Active}(\mathbf{x}^\star)\right\}
  \]
  is linearly independent.

  \item \textbf{Slater's condition.}
  For a convex constrained problem, $\exists\, \bar{\mathbf{x}}\in F$ such
  that
  \[
    h_j(\bar{\mathbf{x}})=0 \quad \forall j,
    \qquad
    g_i(\bar{\mathbf{x}})<0 \quad \forall i.
  \]

  \item \textbf{Weak Slater's condition.}
  For a convex constrained problem, $\exists\, \bar{\mathbf{x}}\in F$ such
  that
  \[
    h_j(\bar{\mathbf{x}})=0 \quad \forall j,
  \]
  and for the inequalities,
  \[
    g_i(\bar{\mathbf{x}})\le 0 \text{ for affine } g_i,
    \qquad
    g_i(\bar{\mathbf{x}})<0 \text{ for the remaining } g_i.
  \]
\end{enumerate}
\end{definition}

The convexity assumptions in the Slater-type conditions refer back to the
material of Chapter~4: in the standard convex setting 
the functions $g_i$ are convex and the equalities $h_j$ are affine.

\subsection{Case 2: one inequality constraint}
Now consider
\[
  f(\mathbf{x}) \to \min_{\mathbf{x}},
  \qquad
  g(\mathbf{x}) \le 0.
\]
If $\mathbf{x}^\star\in \operatorname{int} F$, we are back in the unconstrained
case and $\nabla f(\mathbf{x}^\star)=\mathbf{0}$. So we are more interested in
a boundary point with
$g(\mathbf{x}^\star)=0$. We know that $\nabla g(\mathbf{x})$ will be orthogonal
to the level set and directed outward because, inside the blue region
(see \Cref{fig:ch7-kkt-condition}), the value of $g$ is negative, on the
boundary it is zero, and the gradient is always directed toward increasing
values of the function.
Then see the right panel of
\Cref{fig:ch7-kkt-condition}:
\begin{equation}
  \label{eq:ch7-tangent-one-inequality}
  T_F(\mathbf{x}^\star)
  =
  \left\{
    \mathbf{d}\in\mathbb{R}^n:
    \nabla g(\mathbf{x}^\star)^\top \mathbf{d}\le 0
  \right\}.
\end{equation}

The tangent-cone condition
\[
  \nabla f(\mathbf{x}^\star)^\top \mathbf{d}\ge 0
  \qquad \forall \mathbf{d}\in T_F(\mathbf{x}^\star)
\]
requires $\nabla f(\mathbf{x}^\star)$ to have only one possible direction, like
the one shown in the right panel of \Cref{fig:ch7-kkt-condition}. Equivalently,
\(\exists\, \lambda\ge 0\) such that
\[
  \nabla f(\mathbf{x}^\star)+\lambda \nabla g(\mathbf{x}^\star)=\mathbf{0}.
\]
This is the one-constraint inequality version of KKT. As in the case 1 the same notes about 
constraint qualifications can be made here.

The fourth KKT condition,
\[
  \lambda g(\mathbf{x}^\star)=0,
\]
just covers all possible cases, which can be separated into two cases:
\begin{itemize}
  \item if $\lambda=0$, then $(\mathbf{x}^\star \in \operatorname{int} F$ if $g(\mathbf{x}^\star) <0$ or
  $\mathbf{x}^\star \in \partial F$ if $g(\mathbf{x}^\star) =0$) and
  $\nabla f(\mathbf{x}^\star)=\mathbf{0}$;
  \item if $g(\mathbf{x}^\star)=0$, then
  $\mathbf{x}^\star \in \partial F$.
\end{itemize}
Therefore, it forbids cases when $\lambda >0$  and $g(\mathbf{x}^\star) <0$, so that if we have
inner point $\mathbf{x}^{\star}$, $\nabla f(\mathbf{x}^{\star })$ must be always zero.

\subsection{Case 3: two inequality constraints}
Now consider the next case,
\[
  f(\mathbf{x}) \to \min_{\mathbf{x}},
  \qquad
  g_1(\mathbf{x})\le 0,
  \qquad
  g_2(\mathbf{x})\le 0,
\]
and suppose that both constraints are active at $\mathbf{x}^\star$:
\[
  g_1(\mathbf{x}^\star)=0,
  \qquad
  g_2(\mathbf{x}^\star)=0.
\]
Let $F_i:=\{\mathbf{x}:g_i(\mathbf{x})\le0\}$ for $i=1,2$. The two individual
linearized tangent cones are
\[
  T_{F_i}(\mathbf{x}^\star)
  =
  \left\{\mathbf{d}\in\mathbb{R}^n:
  \nabla g_i(\mathbf{x}^\star)^\top\mathbf{d}\le0\right\},
  \qquad i=1,2.
\]
Since the feasible set is $F=F_1\cap F_2$, a direction is feasible to first
order only if it satisfies both linearized inequalities. Thus
\[
  T_F(\mathbf{x}^\star)
  =
  T_{F_1}(\mathbf{x}^\star)\cap T_{F_2}(\mathbf{x}^\star)
  =
  \left\{\mathbf{d}\in\mathbb{R}^n:
  \nabla g_1(\mathbf{x}^\star)^\top\mathbf{d}\le0,\ 
  \nabla g_2(\mathbf{x}^\star)^\top\mathbf{d}\le0
  \right\}.
\]

\begin{figure}[ht]
    \centering
    \incfig{kkt-general-case}
    \caption[Two active inequalities and the normal cone]{Two active
    inequalities. Each active constraint gives one tangent half-space,
    $T_{F_i}(\mathbf{x}^\star)=\{\mathbf d:
    \nabla g_i(\mathbf{x}^\star)^\top\mathbf d\le0\}$, and the final tangent
    cone is their intersection
    $T_F(\mathbf{x}^\star)=T_{F_1}(\mathbf{x}^\star)\cap
    T_{F_2}(\mathbf{x}^\star)$. The grey cone is the set of vectors having
    nonnegative inner product with every direction in this intersection:
    $\operatorname{cone}\{-\nabla g_1(\mathbf{x}^\star),
    -\nabla g_2(\mathbf{x}^\star)\}$.
    By \Cref{clm:ch7-tangent-necessary}, $\nabla f(\mathbf{x}^\star)$ must lie
    in this grey cone; outside it, the inequality
    $\nabla f(\mathbf{x}^\star)^\top\mathbf d\ge0 \ \forall
    \mathbf d\in T_F(\mathbf{x}^\star)$ fails.}
    \label{fig:ch7-kkt-two-inequalities}
\end{figure}

The tangent-cone condition
\[
  \nabla f(\mathbf{x}^\star)^\top\mathbf{d}\ge0
  \qquad \forall \mathbf{d}\in T_F(\mathbf{x}^\star)
\]
therefore restricts $\nabla f(\mathbf{x}^\star)$ to the grey cone in
\Cref{fig:ch7-kkt-two-inequalities}. In other words,
\[
  \exists\, \lambda_1,\lambda_2\ge0:
  \qquad
  \nabla f(\mathbf{x}^\star)
  =
  -\lambda_1\nabla g_1(\mathbf{x}^\star)
  -\lambda_2\nabla g_2(\mathbf{x}^\star).
\]
Equivalently,
\[
  \nabla f(\mathbf{x}^\star)
  +\lambda_1\nabla g_1(\mathbf{x}^\star)
  +\lambda_2\nabla g_2(\mathbf{x}^\star)
  =
  \mathbf{0},
  \qquad
  \lambda_1,\lambda_2\ge0.
\]

\subsection{The general case and Farkas' lemma}
\label{sec:ch7-farkas}
Once a constraint qualification holds, the tangent cone at a feasible point
$\mathbf{x}^\star$ admits the linear description suggested by the previous cases:
\begin{equation}
  \label{eq:ch7-general-tangent-cone}
  T_F(\mathbf{x}^\star)
  =
  \left\{
    \mathbf{d}\in\mathbb{R}^n:
    \begin{array}{l}
      \nabla g_i(\mathbf{x}^\star)^\top \mathbf{d}\le 0
      \quad \forall i\in \operatorname{Active}(\mathbf{x}^\star)\\
      \nabla h_j(\mathbf{x}^\star)^\top \mathbf{d}=0
      \quad \forall j
    \end{array}
  \right\}.
\end{equation}
So the first-order condition
\eqref{eq:ch7-tangent-necessary} becomes a statement about one vector,
$\nabla f(\mathbf{x}^\star)$, against a cone described by finitely many linear
inequalities and equalities. This is exactly the setting where the cone version
of Farkas' lemma applies.

\begin{lemma}[Farkas]
\label{lem:ch7-farkas}
Let
\[
  K := \left\{B\mathbf{y}+C\mathbf{w}: \mathbf{y}\ge \mathbf{0}\right\},
\]
where $B\in\mathbb{R}^{n\times m}$, $C\in\mathbb{R}^{n\times p}$,
$\mathbf{y}\in\mathbb{R}^m$, and $\mathbf{w}\in\mathbb{R}^p$. Then for
every $\mathbf{g}\in\mathbb{R}^n$, exactly one of the following alternatives
holds:
\begin{enumerate}
  \item $\mathbf{g}\in K$ or
  \item  $\exists\, \mathbf{d}\in\mathbb{R}^n$ such that
  \[
    \mathbf{g}^\top \mathbf{d}<0,
    \qquad
    B^\top \mathbf{d}\ge \mathbf{0},
    \qquad
    C^\top \mathbf{d}=\mathbf{0}.
  \]
\end{enumerate}
\end{lemma}

Geometrically, this lemma is a separation statement for a
point and a convex cone: either the point already lies in the cone, or one can
separate them by a direction.

To use it, introduce the cone generated by the active constraint normals:
\begin{equation}
  \label{eq:ch7-normal-cone}
  N(\mathbf{x}^\star)
  :=
  \left\{
    -\sum_{i\in\operatorname{Active}(\mathbf{x}^\star)} \lambda_i \nabla g_i(\mathbf{x}^\star)
    -\sum_{j=1}^p \mu_j \nabla h_j(\mathbf{x}^\star):
    \lambda_i\ge 0
  \right\}.
\end{equation}
Applying \Cref{lem:ch7-farkas} to
$\mathbf{g}=\nabla f(\mathbf{x}^\star)$ and the linear system
\eqref{eq:ch7-general-tangent-cone} gives two possibilities:
\begin{enumerate}
  \item $\nabla f(\mathbf{x}^\star)\in N(\mathbf{x}^\star)$ or
  \item $\exists\, \mathbf{d}\in\mathbb{R}^n$ such that
  \[
    \begin{aligned}
      \nabla f(\mathbf{x}^\star)^\top \mathbf{d} &<0, \\
      \nabla g_i(\mathbf{x}^\star)^\top \mathbf{d} &\le 0
      \quad \forall i \in\operatorname{Active}(\mathbf{x}^\star), \\
      \nabla h_j(\mathbf{x}^\star)^\top \mathbf{d} &=0
      \quad \forall j.
    \end{aligned}
  \]
\end{enumerate}
But the second alternative gives $\mathbf{d}\in T_F(\mathbf{x}^\star)$ with
$\nabla f(\mathbf{x}^\star)^\top\mathbf{d}<0$, contradicting
\Cref{clm:ch7-tangent-necessary}. Hence only the first alternative can occur.

This finishes a sketch of the proof for \cref{thm:ch7-kkt-first}.


\section{Second-order KKT conditions}
\label{sec:ch7-second-order}
If the stronger inequality
\[
  \nabla f(\mathbf{x}^\star)^\top \mathbf{d}>0
  \qquad \forall \mathbf{d}\in T_F(\mathbf{x}^\star)\setminus\{\mathbf{0}\}
\]
holds, then every nontrivial feasible first-order motion increases the
objective, so $\mathbf{x}^\star$ is already a local minimizer. The subtle case is
when the first-order term vanishes in some feasible directions. We then
pass to a second-order test.

Assume that first-order KKT holds at $(\mathbf{x}^\star,\boldsymbol{\lambda},
\boldsymbol{\mu})$. For every $\mathbf{d}\in T_F(\mathbf{x}^\star)$,
\[
  \nabla f(\mathbf{x}^\star)^\top \mathbf{d}
  = -\sum_{i=1}^m \lambda_i \nabla g_i(\mathbf{x}^\star)^\top \mathbf{d}
  - \sum_{j=1}^p \mu_j \nabla h_j(\mathbf{x}^\star)^\top \mathbf{d}.
\]
Since $\mathbf{d}\in T_F(\mathbf{x}^\star)$, the equality terms vanish and the
inequality terms are nonnegative after multiplication by $-\lambda_i$:
\[
  -\lambda_i \nabla g_i(\mathbf{x}^\star)^\top \mathbf{d}\ge 0.
\]
Therefore a direction with $\nabla f(\mathbf{x}^\star)^\top \mathbf{d} = 0$
must satisfy
\[
  \nabla g_i(\mathbf{x}^\star)^\top \mathbf{d}=0
  \qquad \forall i \in\operatorname{Active}(\mathbf{x}^\star)\text{ with } \lambda_i>0.
\]
This leads to the critical cone.

\begin{definition}[Critical cone]
\label{def:ch7-critical-cone}
Given a feasible point $\mathbf{x}^\star$ and multipliers
$\boldsymbol{\lambda}\ge \mathbf{0}$, the critical cone is
\begin{equation}
  \label{eq:ch7-critical-cone}
  C(\mathbf{x}^\star,\boldsymbol{\lambda})
  :=
  \left\{
    \mathbf{d}\in T_F(\mathbf{x}^\star):
    \begin{array}{l}
      \nabla g_i(\mathbf{x}^\star)^\top \mathbf{d}=0 
      \quad \forall i\in\operatorname{Active}(\mathbf{x}^\star) 
      \text{ and } \lambda_i>0
    \end{array}
  \right\}.
\end{equation}
\end{definition}

\begin{theorem}[Second-order KKT necessity]
\label{thm:ch7-kkt-second-necessary}
Assume that $f,g_i,h_j\in C^2$, let $\mathbf{x}^\star$ be a local minimizer of
\eqref{eq:ch7-problem}, and assume that a constraint qualification holds at
$\mathbf{x}^\star$. Then \(\exists\, (\boldsymbol{\lambda}^\star,\boldsymbol{\mu}^\star)\)
satisfying the
first-order KKT conditions from \Cref{thm:ch7-kkt-first} and, in addition,
\begin{equation}
  \label{eq:ch7-second-order-necessary}
  \mathbf{d}^\top
  \nabla_{\mathbf{x}\mathbf{x}}^2
  L(\mathbf{x}^\star,\boldsymbol{\lambda}^\star,\boldsymbol{\mu}^\star)
  \mathbf{d}
  \ge 0
  \qquad
  \forall \mathbf{d}\in C(\mathbf{x}^\star,\boldsymbol{\lambda}^\star).
\end{equation}
\end{theorem}

\begin{theorem}[Second-order KKT sufficiency]
\label{thm:ch7-kkt-second-sufficient}
Assume that $f,g_i,h_j\in C^2$, that a constraint qualification holds at the
feasible point $\mathbf{x}^\star$, and that
$(\boldsymbol{\lambda}^\star,\boldsymbol{\mu}^\star)$ satisfies the first-order
KKT conditions. If, in addition,
\begin{equation}
  \label{eq:ch7-second-order-sufficient}
  \mathbf{d}^\top
  \nabla_{\mathbf{x}\mathbf{x}}^2
  L(\mathbf{x}^\star,\boldsymbol{\lambda}^\star,\boldsymbol{\mu}^\star)
  \mathbf{d}
  > 0
  \qquad
  \forall \mathbf{d}\in C(\mathbf{x}^\star,\boldsymbol{\lambda}^\star)
  \setminus \{\mathbf{0}\},
\end{equation}
then $\mathbf{x}^\star$ is a local minimizer.
\end{theorem}

The point of \eqref{eq:ch7-second-order-necessary} and
\eqref{eq:ch7-second-order-sufficient} is that the Hessian of the Lagrangian is
not checked on the whole space $\mathbb{R}^n$, but only on the critical
directions left undecided by first-order KKT. However, sometimes we check the
Hessian inequality on the whole space, because there are cases when it is hard
to explicitly construct the critical cone.

\section{Worked examples}
\label{sec:ch7-examples}
\subsection{A closed half-line}
\begin{example}[Minimizing on $x\ge 0$]
\label{ex:ch7-half-line}
Consider
\[
  \min_x x^2
  \qquad \text{subject to } x\ge 0.
\]
To fit the standard form $g(x)\le 0$, write the constraint as $-x\le 0$.
Then
\[
  L(x,\lambda)=x^2-\lambda x.
\]
The KKT system is
\[
  2x-\lambda=0,
  \qquad
  x\ge 0,
  \qquad
  \lambda\ge 0,
  \qquad
  \lambda x=0.
\]
There are two cases.
\begin{itemize}
  \item If $\lambda=0$, then $2x=0$, so $x=0$.
  \item If $\lambda>0$, then complementary slackness gives $x=0$, and then
  $2x-\lambda=0$ forces $\lambda=0$, a contradiction.
\end{itemize}
So the only KKT point is $(x^\star,\lambda^\star)=(0,0)$. Moreover,
\[
  \nabla_{xx}^2 L(x,\lambda)=2>0,
\]
so the second-order sufficient condition holds and $x^\star=0$ is a local
minimum.
\end{example}

\subsection{An open feasible set}
\begin{example}[The infimum is not attained on $x>0$]
\label{ex:ch7-open-set}
Consider
\[
  \min_x x^2
  \qquad \text{subject to } x>0.
\]
The feasible set is open, so every feasible point is an interior point. Hence
the necessary condition for a local minimizer is the unconstrained one:
\[
  \frac{d}{dx}(x^2)=2x=0.
\]
Its only solution is $x=0$, but $x=0$ is infeasible. Therefore the problem has
no minimizer. The infimum equals $0$, but it is approached only as
$x\downarrow 0$.
\end{example}

\subsection{A true minimizer where KKT fails because the constraint qualification fails}
\begin{example}[Why a constraint qualification matters]
\label{ex:ch7-regularity-failure}
Consider
\[
  \min_x x
  \qquad \text{subject to } x^2\le 0.
\]
The feasible set consists of the single point $x=0$, so $x^\star=0$ is the
unique minimizer. Now test KKT. The Lagrangian is
\[
  L(x,\lambda)=x+\lambda x^2.
\]
The KKT system becomes
\[
  1+2\lambda x=0,
  \qquad
  x^2\le 0,
  \qquad
  \lambda\ge 0,
  \qquad
  \lambda x^2=0.
\]
Again split into two cases.
\begin{itemize}
  \item If $\lambda=0$, then the stationarity equation becomes $1=0$, which is
  impossible.
  \item If $\lambda>0$, then complementary slackness yields $x=0$, and the
  stationarity equation again becomes $1=0$.
\end{itemize}
So KKT has no solution, even though the problem has a minimizer. This is not a
contradiction: the constraint is not regular at the minimizer because
\[
  \nabla(x^2)\big|_{x=0}=0.
\]
This is exactly the kind of bad reformulation warned about at the beginning of
the chapter.
\end{example}

\subsection{A quadratic on the unit disk}
\begin{example}[Second-order KKT distinguishes true minima from a false one]
\label{ex:ch7-disk}
Consider
\[
  \min_{x,y} \ x^2-y^2
  \qquad \text{subject to } x^2+y^2\le 1.
\]
Slater's condition holds: for example, $(0,0)$ is strictly feasible since
$0^2+0^2-1<0$. Therefore the KKT conditions are applicable.
The Lagrangian is
\[
  L(x,y,\lambda)=x^2-y^2+\lambda(x^2+y^2-1).
\]
The KKT system is
\begin{equation}
  \label{eq:ch7-disk-kkt}
  \begin{aligned}
    2x+2\lambda x &= 2x(1+\lambda)=0, \\
    -2y+2\lambda y &= 2y(\lambda-1)=0, \\
    x^2+y^2 &\le 1, \\
    \lambda &\ge 0, \\
    \lambda(x^2+y^2-1)&=0.
  \end{aligned}
\end{equation}

First solve \eqref{eq:ch7-disk-kkt}.
\begin{itemize}
  \item If $\lambda=0$, then $x=0$ and $y=0$, so $(0,0)$ satisfies KKT.
  \item If $\lambda>0$, then complementary slackness gives $x^2+y^2=1$. The
  first two equations imply $x=0$ and $\lambda=1$, hence $y=\pm 1$. So
  $(0,1)$ and $(0,-1)$ are also KKT points.
\end{itemize}

Now use the second-order test. The Hessian of the Lagrangian is
\[
  \nabla_{\mathbf{x}\mathbf{x}}^2 L(x,y,\lambda)
  =
  \begin{bmatrix}
    2(1+\lambda) & 0 \\
    0 & 2(\lambda-1)
  \end{bmatrix}.
\]

At $(0,0)$ we have $\lambda=0$, so
\[
  \nabla_{\mathbf{x}\mathbf{x}}^2 L(0,0,0)
  =
  \begin{bmatrix}
    2 & 0 \\
    0 & -2
  \end{bmatrix}.
\]
Here the constraint is inactive, so the tangent cone is all of $\mathbb{R}^2$,
and the critical cone is also $\mathbb{R}^2$. The Hessian is indefinite on
$\mathbb{R}^2$, therefore $(0,0)$ is \emph{not} a local minimum.

At $(0,1)$ we have $\lambda=1$, hence
\[
  \nabla_{\mathbf{x}\mathbf{x}}^2 L(0,1,1)
  =
  \begin{bmatrix}
    4 & 0 \\
    0 & 0
  \end{bmatrix}.
\]
The active constraint gradient is
\[
  \nabla g(0,1)=
  \nabla(x^2+y^2-1)\big|_{(0,1)}
  =
  \begin{bmatrix}
    0 \\
    2
  \end{bmatrix},
\]
so
\[
  T_F(0,1)
  = \left\{
    \mathbf{d}=(d_1,d_2)^\top:
    0\cdot d_1 + 2d_2 \le 0
  \right\}.
\]
Because $\lambda=1>0$, the critical cone imposes equality:
\[
  C((0,1),1)
  = \left\{
    \mathbf{d}\in T_F(0,1): 2d_2=0
  \right\}
  = \left\{
    \mathbf{d}: d_2=0
  \right\}.
\]
Therefore, \(\forall \mathbf{d}\in C((0,1),1)\setminus\{\mathbf{0}\}\),
\[
  \mathbf{d}^\top
  \nabla_{\mathbf{x}\mathbf{x}}^2 L(0,1,1)
  \mathbf{d}
  = 4d_1^2 > 0.
\]
So $(0,1)$ is a local minimum.

Exactly the same argument applies to $(0,-1)$, where
\[
  T_F(0,-1)
  = \left\{
    \mathbf{d}=(d_1,d_2)^\top:
    -2d_2 \le 0
  \right\},
  \qquad
  C((0,-1),1)=\{\mathbf{d}: d_2=0\},
\]
and again the quadratic form equals $4d_1^2>0$ on nonzero critical
directions. Hence $(0,-1)$ is also a local minimum.
\end{example}

\subsection{A determinant maximization example}
\begin{example}[Maximizing $\det(\mathbf{X})$ under one trace constraint]
\label{ex:ch7-det-max}
Recall from Chapter~4 that $\mathbb{S}_{++}^n$ denotes the cone of symmetric
positive definite matrices. Let $\mathbf{A}\in \mathbb{S}_{++}^n$ and $b>0$,
and consider
\[
  \max_{\mathbf{X}\in \mathbb{S}_{++}^n}
  \det(\mathbf{X})
  \qquad
  \text{subject to }
  \langle \mathbf{A},\mathbf{X}\rangle \le b,
\]
where
\[
  \langle \mathbf{A},\mathbf{X}\rangle
  := \operatorname{tr}(\mathbf{A}^\top \mathbf{X}).
\]
To fit the chapter's minimization form, rewrite the problem as
\[
  \min_{\mathbf{X}\in \mathbb{S}_{++}^n} -\det(\mathbf{X})
  \qquad
  \text{subject to }
  \operatorname{tr}(\mathbf{A}^\top \mathbf{X})-b \le 0.
\]
The constraint qualification holds because the inequality constraint is affine
in $\mathbf{X}$.
The Lagrangian is
\[
  L(\mathbf{X},\lambda)
  =
  -\det(\mathbf{X})
  +\lambda\bigl(\operatorname{tr}(\mathbf{A}^\top \mathbf{X})-b\bigr).
\]

For $\mathbf{X}\in \mathbb{S}_{++}^n$, the differential of the determinant is
\[
  d(\det(\mathbf{X}))
  =
  \det(\mathbf{X}) \operatorname{tr}(\mathbf{X}^{-1} d\mathbf{X}),
\]
so
\begin{align*}
  dL
  &=
  -\det(\mathbf{X}) \operatorname{tr}(\mathbf{X}^{-1} d\mathbf{X})
  + \lambda \operatorname{tr}(\mathbf{A}^\top d\mathbf{X}) \\
  &=
  \operatorname{tr}\Bigl(
    \bigl(-\det(\mathbf{X})\mathbf{X}^{-1} + \lambda \mathbf{A}\bigr)
    d\mathbf{X}
  \Bigr),
\end{align*}
because $\mathbf{X}$ and $\mathbf{A}$ are symmetric. Hence the KKT system is
\begin{equation}
  \label{eq:ch7-det-kkt}
  \begin{aligned}
    -\det(\mathbf{X})\mathbf{X}^{-1} + \lambda \mathbf{A} &= \mathbf{0}, \\
    \operatorname{tr}(\mathbf{A}^\top \mathbf{X}) &\le b, \\
    \lambda &\ge 0, \\
    \lambda\bigl(\operatorname{tr}(\mathbf{A}^\top \mathbf{X})-b\bigr) &= 0.
  \end{aligned}
\end{equation}

Now solve \eqref{eq:ch7-det-kkt}. If $\lambda=0$, then the first equation gives
\[
  \det(\mathbf{X})\mathbf{X}^{-1}=\mathbf{0},
\]
which is impossible because $\mathbf{X}\succ 0$ implies
$\det(\mathbf{X})>0$ and $\mathbf{X}$ is invertible. Therefore $\lambda>0$, and
complementary slackness forces the constraint to be active:
\[
  \operatorname{tr}(\mathbf{A}^\top \mathbf{X})=b.
\]
From stationarity we obtain
\[
  \mathbf{X}
  =
  \frac{\det(\mathbf{X})}{\lambda}\mathbf{A}^{-1}
  =:\alpha \mathbf{A}^{-1}.
\]
Substituting this into the active constraint gives
\[
  b
  =
  \operatorname{tr}(\mathbf{A}^\top \mathbf{X})
  =
  \alpha \operatorname{tr}(\mathbf{A}^\top \mathbf{A}^{-1})
  =
  \alpha \operatorname{tr}(\mathbf{I})
  =
  \alpha n.
\]
Hence
\[
  \alpha = \frac{b}{n},
  \qquad
  \mathbf{X}^\star = \frac{b}{n}\mathbf{A}^{-1}.
\]

This matrix is feasible and lies in $\mathbb{S}_{++}^n$, so the optimizer is
\[
  \boxed{
    \mathbf{X}^\star = \frac{b}{n}\mathbf{A}^{-1}
  }.
\]
\end{example}

\subsection{DeepFool as a constrained quadratic subproblem}
\begin{example}[DeepFool adversarial attack]
\label{ex:ch7-deepfool}
This example shows how the KKT machinery of this chapter also appears in a
modern machine-learning application. Let
\[
  f:\mathbb{R}^n \to \mathbb{R}^K,
  \qquad
  \hat k(\mathbf{x}) := \arg\max_{1\le k\le K} f_k(\mathbf{x}),
\]
where $f_k(\mathbf{x})$ is the score of class $k$. For an input
$\mathbf{x}$ with predicted class
\[
  y := \hat k(\mathbf{x}),
\]
the DeepFool attack seeks a small perturbation $\mathbf{r}$ that changes the
predicted label. The basic objective is
\[
  \min_{\mathbf{r}} \frac12 \|\mathbf{r}\|_2^2
  \qquad
  \text{subject to }
  \hat k(\mathbf{x}+\mathbf{r}) \neq y.
\]

This constraint is nonlinear, so DeepFool replaces it by a first-order model.
Linearizing each score around the current point gives
\[
  f_k(\mathbf{x}+\mathbf{r})
  \approx
  f_k(\mathbf{x}) + \nabla f_k(\mathbf{x})^\top \mathbf{r}.
\]
The predicted class changes when at least one competing score catches the
current winning score. Locally, we therefore ask how far one must move before
the first-order model of $f_j$ ties the first-order model of $f_y$.
For a fixed competing class $j\neq y$, this linearized boundary between class
$y$ and class $j$ is
\[
  f_y(\mathbf{x}) + \nabla f_y(\mathbf{x})^\top \mathbf{r}
  =
  f_j(\mathbf{x}) + \nabla f_j(\mathbf{x})^\top \mathbf{r}.
\]
So the local step is the equality-constrained problem
\begin{equation}
  \label{eq:ch7-deepfool-subproblem}
  \min_{\mathbf{r}} \frac12 \|\mathbf{r}\|_2^2
  \qquad
  \text{subject to }
  (\nabla f_y(\mathbf{x})-\nabla f_j(\mathbf{x}))^\top \mathbf{r}
  + f_y(\mathbf{x}) - f_j(\mathbf{x}) = 0.
\end{equation}

This is exactly the kind of problem treated earlier in the chapter: a strictly
convex quadratic objective with one affine equality constraint. Its Lagrangian
is
\[
  L(\mathbf{r},\mu)
  =
  \frac12 \|\mathbf{r}\|_2^2
  +
  \mu\Bigl(
    (\nabla f_y(\mathbf{x})-\nabla f_j(\mathbf{x}))^\top \mathbf{r}
    + f_y(\mathbf{x}) - f_j(\mathbf{x})
  \Bigr).
\]
The KKT conditions are
\[
  \nabla_{\mathbf{r}}L(\mathbf{r},\mu)
  =
  \mathbf{r} + \mu(\nabla f_y(\mathbf{x})-\nabla f_j(\mathbf{x}))
  = \mathbf{0},
\]
together with the equality constraint from
\eqref{eq:ch7-deepfool-subproblem}. From stationarity,
\[
  \mathbf{r}
  =
  -\mu(\nabla f_y(\mathbf{x})-\nabla f_j(\mathbf{x})).
\]
Substituting this into the constraint yields
\[
  -\mu\|\nabla f_y(\mathbf{x})-\nabla f_j(\mathbf{x})\|_2^2
  + f_y(\mathbf{x}) - f_j(\mathbf{x})
  = 0,
\]
so
\[
  \mu
  =
  \frac{f_y(\mathbf{x})-f_j(\mathbf{x})}
  {\|\nabla f_y(\mathbf{x})-\nabla f_j(\mathbf{x})\|_2^2}.
\]
Therefore the smallest-norm linearized perturbation that reaches class
$j$ is
\begin{equation}
  \label{eq:ch7-deepfool-rj}
  \mathbf{r}_j
  =
  -
  \frac{f_j(\mathbf{x})-f_y(\mathbf{x})}
  {\|\nabla f_j(\mathbf{x})-\nabla f_y(\mathbf{x})\|_2^2}
  \bigl(\nabla f_j(\mathbf{x})-\nabla f_y(\mathbf{x})\bigr).
\end{equation}
Its norm is
\[
  \|\mathbf{r}_j\|_2
  =
  \frac{|f_j(\mathbf{x})-f_y(\mathbf{x})|}
  {\|\nabla f_j(\mathbf{x})-\nabla f_y(\mathbf{x})\|_2}.
\]
Hence the locally closest class boundary is obtained by choosing
\begin{equation}
  \label{eq:ch7-deepfool-class-choice}
  \ell(\mathbf{x})
  :=
  \arg\min_{j\neq y}
  \frac{|f_j(\mathbf{x})-f_y(\mathbf{x})|}
  {\|\nabla f_j(\mathbf{x})-\nabla f_y(\mathbf{x})\|_2}.
\end{equation}
The corresponding step is
\[
  \mathbf{r}
  =
  \mathbf{r}_{\ell(\mathbf{x})}.
\]

DeepFool then applies this construction iteratively. Starting from
$\mathbf{x}_0=\mathbf{x}$, it keeps updating
\[
  \mathbf{x}_{i+1} = \mathbf{x}_i + \mathbf{r}_i
\]
as long as the predicted class remains equal to the original label $y$.
At step $i$, one recomputes the gradients and scores at $\mathbf{x}_i$,
selects
\[
  \ell_i
  =
  \arg\min_{j\neq y}
  \frac{|f_j(\mathbf{x}_i)-f_y(\mathbf{x}_i)|}
  {\|\nabla f_j(\mathbf{x}_i)-\nabla f_y(\mathbf{x}_i)\|_2},
\]
and uses
\[
  \mathbf{r}_i
  =
  -
  \frac{f_{\ell_i}(\mathbf{x}_i)-f_y(\mathbf{x}_i)}
  {\|\nabla f_{\ell_i}(\mathbf{x}_i)-\nabla f_y(\mathbf{x}_i)\|_2^2}
  \bigl(\nabla f_{\ell_i}(\mathbf{x}_i)-\nabla f_y(\mathbf{x}_i)\bigr).
\]
The final output is the accumulated perturbation, often with a small
overshoot factor:
\[
  \hat{\mathbf{r}}
  =
  (1+\eta)(\mathbf{x}_m-\mathbf{x}_0),
  \qquad
  \eta \approx 0.02,
\]
where $m$ is the first index for which
$\hat k(\mathbf{x}_m)\neq y$.
The overshoot is a practical correction: each computed step is the shortest
step to a locally linearized boundary, so without the factor the final point may
lie exactly on the approximate boundary or may still remain on the original side
of the true nonlinear boundary. Multiplying by $1+\eta$ moves the point slightly
farther in the same direction and makes the label change robust to curvature and
numerical instabilities.

\end{example}

